The study of citation (impact) relationships and citation metrics has continued to advance from a wide range of perspectives in recent years. Li et al. \cite{Li1} analyzed the long term characteristics of citations for 23 million papers and elucidated their group structure. In addition, Yu et al. \cite{Yu1} performed a structural analysis of the 100 most cited papers in the field of social psychology, revealing the influence structures of major publication years, international collaborations, and journals, thereby providing insights into the research network structure of cited papers.

Focusing on the journal level, Leydesdorff \cite{Ley1,Ley2,Ley3,Ley4} constructed and examined maps of inter journal relations. Amano et al. \cite{Ama1,Ama2} carried out large scale clustering of science, technorogy, engneering and mathematical journals and clarified citation relationships across fields. These studies attempt to capture, from a macro perspective, the conditions under which highly cited papers or high impact journals arise and the pathways of their influence.

At the same time, the systematization of structural factors that contribute to citation at the paper level is progressing, suggesting that multiple factors such as collaborative research forms, publication strategies, and allocation of research resources may affect citation counts. Hill et al. \cite{Hil1} pointed out that a researcher's strategy of shifting research topics may actually reduce subsequent impact. Lin et al. \cite{Lin1}, using fulltext data of cited papers, observed that the location, context, and similarity to the citer change over time, indicating a transformation of the role of citations that cannot be captured by citation counts alone. Kousha et al. \cite{Ko1} investigated how the "readability" of abstracts and the "quality" of papers in the information science field influence citation numbers.  

The studies introduced above primarily elucidate various attributes related to citability, but research that multilayeredly captures the long term structural characteristics common to highly cited papers is limited. This study provides a perspective different from existing research by integratively analyzing elements such as citation lifespan, journal structure, and thematic structure.
